We list the boundaries of the result plainly, because several of them are load-bearing.

\paragraph{Memory is $O(d^2)$, not $O(d)$.}
The $O(dN)$ claim is about \emph{time}. The method stores a $d\times d$ inverse curvature
estimate and two $d\times d$ covariance accumulators, so memory is $O(d^2)$; this is
inherited from the independent-data method we extend and is not something we improve. For
$d$ in the thousands this is the binding constraint, and a sketched or diagonal-plus-low-rank
curvature representation---as in \citet{kuang2026online,wang2026inference,godichonbaggioni2026complexity}---%
would be the natural remedy. Whether the block-gradient long-run covariance estimator
survives such a compression is open.

\paragraph{Geometric mixing is a real restriction, and the block length has to respect it.}
\Cref{thm:lrv} buys $b\asymp\log N$ from the geometric decay in
\Cref{lem:davydov}. Under algebraic mixing the flat-top bias is polynomial rather than
exponential and the admissible $b$ grows accordingly; long-memory streams are outside the
theory altogether. This is not merely formal: our air-quality stream has a residual lag-one
autocorrelation of $0.92$ and a variance-inflation factor around $20$, and there the
block length needed for calibrated intervals is an order of magnitude larger than for the
traffic stream (\Cref{tab:real}). The free diagnostic
$r_j=|u_j^{\!\top}(\Sft-\Sbm)u_j|/(u_j^{\!\top}\Sft u_j)$ detects the problem---it estimates
the relative block-length bias of plain batch means---but it does not remove it. It could be
turned into an automatic rule by maintaining covariance accumulators at several block scales
$b,2b,4b,\dots$ at once (one extra $O(d^2)$ update per block per scale, all from the same
block gradients) and reporting from the smallest scale at which $r_j$ falls below a
tolerance; we did not implement this, and we do not claim an automatic rule for $b$.

\paragraph{$\Sft$ need not be positive semidefinite.}
This is intrinsic to flat-top and lugsail estimators \citep[Remark~5]{singh2025utility}.
Marginal intervals need only positive scalars $u_j^{\!\top}\Sft u_j$, and the documented
fallback to $\Sbm$ never triggered in any of our runs, but a user who needs the full matrix
(for a joint test, say) must project it.

\paragraph{The curvature form is required.}
Everything rests on
$\nabla^2F(\theta)=\E[\alpha\,\Psi\Psi^{\!\top}]$, which holds for generalised linear
models, ridge and softmax regression, the geometric median and Gauss--Newton problems, but
not for a general objective. Without it there is no rank-one recursion and no
inversion-free update.

\paragraph{Warm-up is a real requirement, not an implementation detail.}
The $1/t$ schedule with a poor preconditioner leaves a slowly decaying component that
dominates the sampling error at realistic $N$: with a negligible warm-up, coverage of
nominal $95\%$ intervals fell to $\CovNoWarmup$ in our autoregressive design even though the
variance estimate stayed accurate (\Cref{tab:ablation}(c)). We give a scaling rule
($c_{\mathrm w}d$ curvature updates, with $c_{\mathrm w}=50$ throughout) and show
insensitivity above it, but this is a tuning constant, and the asymptotic theory is silent
about it. It also has a cost consequence we state rather than bury: the warm-up performs
$c_{\mathrm w}d$ rank-one updates at $O(d^2)$ each, so the total time is
$O(dN+c_{\mathrm w}d^3)$ and the advertised $O(dN)$ only bites once $N\gtrsim c_{\mathrm w}d^2$
(with $c_{\mathrm w}=50$ and $d=80$ that means $N\gtrsim3\times10^5$; at the $N=2\times10^5$ of
\Cref{tab:cost} the warm-up is comparable to the streaming pass).

\paragraph{Rates and what we do not prove.}
We give an almost sure rate and a central limit theorem, not a Berry--Esseen bound; the
first-order Markovian literature has the latter for linear stochastic approximation
\citep{samsonov2025statistical}, and we do not match it. We attain the semiparametric
efficiency bound of \citet{li2023markovian} rather than establishing it. Our
$\widetilde O(N^{-1/2})$ rate for $\Sft$ is not comparable to the minimax rate for
Hessian-free covariance estimation from an iterate path \citep{ni2026online}, because we use
more information (block gradients and a maintained curvature estimate); we make no claim
about optimality within our own information set.

\paragraph{Scope of the empirical evidence.}
All designs are convex $M$-estimation with $d\le160$ and $N\le4\times10^5$; there is no
deep-learning experiment and the theory would not support one. The dependence in the
simulated designs is stationary and geometrically decaying by construction. On the real
streams we do not know the truth, which is why we run two protocols; the semi-synthetic
protocol has an exact target but a synthetic (autoregressive) error process, and the fully
real protocol has a real error process but a target that is itself estimated---we correct
the nominal level for that exactly, but the correction assumes the segments are
approximately exchangeable, which non-stationarity would violate. Neither protocol alone
would be convincing; together they bracket the honest answer.

\paragraph{One thing we deliberately do not claim.}
Gapping improves the curvature estimate at matched cost (\Cref{prop:gap}, \Cref{fig:gap}),
but the curvature estimate enters the limit distribution only through a rate, so the
end-to-end effect on the point estimate and on coverage is second order and we could not
measure it reliably. The case for gapping is that it is free, deterministic, and provably
better than the randomised alternative---not that it changes the headline numbers.
